% 《大学物理简明教程》（第5版）2.3.13--2.3.25
% 来源第53--54页（印刷页与PDF页一致），逐题研究整理。待独立复核。

\subsection{解答题：2.3.13--2.3.25}

\sourceproblem{2.3.13}{53}{53}
\problemtext 质点在$Oxy$面上运动，$\bm r=a\cos\omega t\,\bm i+b\sin\omega t\,\bm j$，$a>b>0$，求作用力、在$A(a,0)$与$B(0,b)$的动能、$A\to B$时作用力做功及该力是否保守。
\workings
逐次微分
\[
\bm v=-a\omega\sin\omega t\,\bm i+b\omega\cos\omega t\,\bm j,
\quad
\bm F=m\ddot{\bm r}=-m\omega^2\bm r.
\]
设$t=0$为$A$，$t=\pi/(2\omega)$为$B$，
\[
\boxed{K_A=\frac12m b^2\omega^2},\qquad
\boxed{K_B=\frac12m a^2\omega^2}.
\]
由动能定理
\[
\boxed{W_{A\to B}=K_B-K_A
=\tfrac12m\omega^2(a^2-b^2)>0}.
\]
因$\bm F=-\nabla U$，其中
$U=\frac12m\omega^2(x^2+y^2)$，
故是\textbf{保守力}。
\crosscheck 力的旋度为零；对$A,B$直接计算$U_A-U_B$与$W$一致。

\sourceproblem{2.3.14}{53}{53}
\problemtext 质点受恒力$\bm F=(7\bm i-6\bm j)\,\mathrm N$，从原点运动到$\bm r=(-3\bm i+4\bm j+16\bm k)\,\mathrm m$。求该力做功、用时$0.6\,\mathrm s$时平均功率、质量$1\,\mathrm{kg}$时动能改变量。
\workings
\[
W=\bm F\cdot\Delta\bm r=7(-3)+(-6)4+0(16)
=\boxed{-45\,\mathrm J}.
\]
平均功率$\boxed{\bar P=W/\Delta t=-75\,\mathrm W}$，
动能定理给$\boxed{\Delta K=-45\,\mathrm J}$。
\crosscheck 动能变化可以为负，但需要初动能至少$45\,\mathrm J$才能到达题设终点；力在$z$方向不做功。

\sourceproblem{2.3.15}{53}{53}
\problemtext 锤击铁钉入木板，木板对铁钉阻力与已进入深度$x$成正比：$F=kx$。相同速度、相同质量的锤两次撞击，第一次打入$1\,\mathrm{cm}$。求第二次还能打入多深。
\workings
设每次传给铁钉的相同有效动能为$E$。第一次从$x=0$进入$x_1$，
$E=\int_0^{x_1}kx\,dx=\frac12kx_1^2$。
第二次从$x_1$推进到$x_2$，
\[
E=\frac12k(x_2^2-x_1^2)
\Longrightarrow x_2^2=2x_1^2.
\]
所以总深度$x_2=\sqrt2 x_1$，
\[
\boxed{\Delta x_2=(\sqrt2-1)x_1
=(\sqrt2-1)\,\mathrm{cm}\approx0.414\,\mathrm{cm}}.
\]
\crosscheck 第一次进入的阻力起点为零，第二次阻力起点为$kx_1$，故第二次打入深度更短。

\sourceproblem{2.3.16}{53}{53}
\problemtext 两个质量可忽略的弹簧$A,B$竖直串联，劲度系数分别为$k_1,k_2$，下端悬质量$M$重物，系统静止时求两弹簧伸长量之和与弹性势能之比。
\workings
串联平衡时两弹簧拉力均为$Mg$，
\[
x_1=\frac{Mg}{k_1},\quad x_2=\frac{Mg}{k_2},\quad
x=x_1+x_2=Mg\left(\frac1{k_1}+\frac1{k_2}\right).
\]
弹性势能
$U=\frac12 k_1x_1^2+\frac12k_2x_2^2
=\frac12(Mg)^2(1/k_1+1/k_2)$，故
\[
\boxed{\frac{x}{U}=\frac{2}{Mg}}.
\]
\crosscheck $x/U$量纲为$\mathrm{m/J=1/N}$，与$2/(Mg)$一致。

\sourceproblem{2.3.17}{54}{54}
\problemtext 太阳系以外的月地引力模型：地球质量$M_E=5.98\times10^{24}\,\mathrm{kg}$，月球质量$M_M=7.35\times10^{22}\,\mathrm{kg}$，地月中心距$D=3.84\times10^8\,\mathrm m$，月半径$R_M=1.74\times10^6\,\mathrm m$。
求地月连线上引力相互抵消点$P$到月球表面的距离，并以无穷远为势能零点求该处$1\,\mathrm{kg}$质点势能。
\workings
设$P$到\emph{月心}距离为$d$，则到地心距离为$D-d$。
相反方向的引力大小相等：
\[
\frac{GM_E}{(D-d)^2}=\frac{GM_M}{d^2}
\Longrightarrow d=\frac{D}{1+\sqrt{M_E/M_M}}.
\]
因此月表至$P$距离
\[
\boxed{h=d-R_M\approx3.6583\times10^7\,\mathrm m}.
\]
以无穷远为引力势能零点，$m=1\,\mathrm{kg}$，
\[
\boxed{U(P)=-Gm\left(\frac{M_E}{D-d}+\frac{M_M}{d}\right)
\approx-1.283\times10^6\,\mathrm J}
\]
（取$G=6.6743\times10^{-11}$ SI）。
\crosscheck 该点处\emph{力}为零，但引力势能不为零；月心到$P$的距离仍远大于月半径。

\sourceproblem{2.3.18}{54}{54}
\problemtext 质量$m$物体在与水平面成$\theta$角的光滑斜面上，由斜面上端弹簧$k$连接，沿斜面向下为$x$正向，原点在弹簧自然长处。物体在平衡点的动能为$E_{k0}$，求平衡点坐标$x_0$及任意弹簧伸长$x$时动能。
\workings
沿斜面合力$mg\sin\theta-kx$，平衡时
\[
\boxed{x_0=\frac{mg\sin\theta}{k}}.
\]
总势能$U(x)=\frac12kx^2-mgx\sin\theta$，
配方得
\[
U(x)=\frac12k(x-x_0)^2-\frac12kx_0^2.
\]
机械能守恒，利用在$x_0$处动能$E_{k0}$，
\[
\boxed{E_k(x)=E_{k0}-\frac12k(x-x_0)^2}.
\]
\crosscheck 可达区域必须满足$E_k(x)\ge0$，故
$|x-x_0|\le\sqrt{2E_{k0}/k}$。

\sourceproblem{2.3.19}{54}{54}
\problemtext 质量$M$的弧形槽在光滑水平面上自由移动，其内轨道为半径$R$的四分之一圆弧。质量$m$小球从槽顶静止释放，沿无摩擦轨道滑到最低端并脱离，求小球相对于槽的速率。
\workings
小球下降高度$R$，系统无水平外力，水平总动量守恒。槽底切线水平，设小球、槽对地速度分别为$v,V$，取同一直线有
\[
mv+MV=0,\qquad
mgR=\frac12mv^2+\frac12MV^2.
\]
令相对速率$u=|v-V|$，由动量守恒$v=\frac{M}{M+m}u$、
$V=-\frac{m}{M+m}u$，代入能量得
\[
mgR=\frac12\frac{mM}{M+m}u^2
\quad\Longrightarrow\quad
\boxed{u=\sqrt{\frac{2gR(M+m)}{M}}}.
\]
\crosscheck $M\to\infty$时槽固定，$u\to\sqrt{2gR}$；有限$M$时槽反冲使相对速率增大。

\sourceproblem{2.3.20}{54}{54}
\problemtext 质量$m$质点在$Oxy$面上的位置为$(x_1,y_1)$，瞬时速度$\bm v=v_x\bm i+v_y\bm j$。受到沿$x$轴负方向、大小为$f$的力，求相对原点的角动量及该力的力矩。
\workings
$\bm r=x_1\bm i+y_1\bm j$，$\bm p=m(v_x\bm i+v_y\bm j)$，所以
\[
\boxed{\bm L_O=\bm r\times\bm p
=m(x_1v_y-y_1v_x)\bm k}.
\]
$\bm F=-f\bm i$，
\[
\boxed{\bm\tau_O=\bm r\times\bm F=f y_1\bm k}.
\]
\crosscheck 由$dL_z/dt=xF_y-yF_x=f y$，与力矩公式一致。

\sourceproblem{2.3.21}{54}{54}
\problemtext 质量$3\,\mathrm{kg}$的质点在$t=0$处$\bm r_0=4\bm i$（m），
$\bm v_0=\bm i+6\bm j$（m/s），之后恒力$\bm F=5\bm j$（N）作用$3$秒。
求动量改变量与对原点角动量改变量。
\workings
\[
\boxed{\Delta\bm p=\bm F\,\Delta t=15\bm j\ \mathrm{kg\,m/s}}.
\]
由于$x(t)=4+t$（m），力矩$\tau_z=5x(t)=5(4+t)$（N\,m），
\[
\boxed{\Delta\bm L=\int_0^3\bm\tau\,dt
=5\left[4t+\tfrac12 t^2\right]_{0}^{3}\bm k
=82.5\bm k\ \mathrm{kg\,m^2/s}}.
\]
\crosscheck 计算初末角动量$\bm r\times m\bm v$也得到$82.5\bm k$的增量。

\sourceproblem{2.3.22}{54}{54}
\problemtext 哈雷彗星绕太阳运行椭圆轨道；近日点距太阳$r_1=8.75\times10^{10}\,\mathrm m$、速率$v_1=5.46\times10^4\,\mathrm{m/s}$，远日点速率$v_2=9.08\times10^2\,\mathrm{m/s}$。求远日点距离$r_2$。
\workings
引力为中心力，两个拱点的速度与径向垂直，故
\[
mr_1v_1=mr_2v_2
\Longrightarrow
\boxed{r_2=r_1\frac{v_1}{v_2}
\approx5.2616\times10^{12}\,\mathrm m}.
\]
\crosscheck 近日点速率较大，远日点距离也较大，$r_2/r_1=v_1/v_2\approx60.13$。

\sourceproblem{2.3.23}{54}{54}
\problemtext 光滑水平桌面中心有小孔，质量$m$的小球以半径$r_0$、角速度$\omega_0$做圆周运动；绳穿孔下挂质量$M_1$的重物时恰可维持平衡。再加挂质量$M_2$且慢慢重新达到平衡时，求小球新的角速度与轨道半径。
\workings
初始绳张力提供向心力：$M_1g=m\omega_0^2 r_0$。
缓慢调整至新平衡时
$(M_1+M_2)g=m\omega^2r$。
关于孔的总外力矩为零，故
\[
mr^2\omega=mr_0^2\omega_0
\quad\Longrightarrow\quad
\omega=\omega_0(r_0/r)^2.
\]
代回平衡式并与初始式相除：
\[
\frac{M_1+M_2}{M_1}=\left(\frac{r_0}{r}\right)^3.
\]
所以
\[
\boxed{r=r_0\left(\frac{M_1}{M_1+M_2}\right)^{1/3}},\qquad
\boxed{\omega=\omega_0\left(\frac{M_1+M_2}{M_1}\right)^{2/3}}.
\]
\crosscheck $M_2\to0$时两式还原；加重物后$r$减小、$\omega$增大。假设准静态、拉力始终径向。

\sourceproblem{2.3.24}{54}{54}
\problemtext 水流水平管中的体积流量$Q=4.0\times10^3\,\mathrm{cm^3/s}$，
粗、细管面积$S_1=40\,\mathrm{cm^2}$、$S_2=10\,\mathrm{cm^2}$。
(1)求两处流速；(2)以水银--水差压计测压，求两处水银面高度差，忽略粘性损失。
\workings
换算$Q=4.0\times10^{-3}\,\mathrm{m^3/s}$，
$S_1=4.0\times10^{-3}\,\mathrm{m^2}$，$S_2=10^{-3}\,\mathrm{m^2}$，由连续性有
\[
\boxed{v_1=Q/S_1=1\,\mathrm{m/s}},\quad
\boxed{v_2=Q/S_2=4\,\mathrm{m/s}}.
\]
同高水平管伯努利方程给
$p_1-p_2=\frac12\rho_w(v_2^2-v_1^2)=7500\,\mathrm{Pa}$。
差压计关系
$\Delta p=(\rho_{\rm Hg}-\rho_w)g\,\Delta h$，
取$\rho_w=1000,\rho_{\rm Hg}=13600\,\mathrm{kg/m^3}$，
\[
\boxed{\Delta h=\frac{7500}{(13600-1000)9.8}
\approx0.0607\,\mathrm m=6.07\,\mathrm{cm}}.
\]
\crosscheck 粗管速度较小、压力较高；差压计结果是正的液面高度差。

\sourceproblem{2.3.25}{54}{54}
\problemtext $20^\circ\mathrm C$时油密度$\rho=0.90\times10^3\,\mathrm{kg/m^3}$；
粗细均匀供油管1处比2处高$5.0\,\mathrm m$，且$p_1$比$p_2$\emph{低}$1.2\times10^3\,\mathrm{Pa}$。
求体积$V=5.0\,\mathrm{m^3}$油从1到2的机械能损耗。
\workings
管径不变且密度固定，设截面平均速度近似相同，动能差抵消。机械能耗散
\[
E_{\rm loss}=V[(p_1-p_2)+\rho g(z_1-z_2)].
\]
代入$p_1-p_2=-1200\,\mathrm{Pa}$、$z_1-z_2=5.0\,\mathrm m$，
\[
\boxed{E_{\rm loss}=5[-1200+900(9.8)(5)]
=2.145\times10^5\,\mathrm J}.
\]
\crosscheck 若将$p_1$“低”误当“高”，将产生数值偏差$2(1200)(5)=12000\,\mathrm J$；这里保持原书差压符号。
